\textbf{UC-INW-01: Przyjęcie Samochodu na plac}

\begin{longtable}{|p{4cm}|p{11cm}|}
\hline
\textbf{Element} & \textbf{Opis} \\
\hline

Identyfikator & UC-INW-01 \\
\hline

Nazwa & Przyjęcie Samochodu na plac \\
\hline

Opis &
Systemowa rejestracja fizycznego auta zjeżdżającego z lawety na plac salonu samochodowego na podstawie numeru VIN i wprowadzenie go do stanu magazynowego. \\
\hline

Aktor główny & Logistyk / Magazynier \\
\hline

Aktorzy pomocniczy & System ERP Importera/Producenta  \\
\hline

Warunki wstępne &
W systemie istnieje awizowana dostawa z fabryki lub utworzony wirtualny "Slot Produkcyjny" przypisany do salonu samochodowego. \\
\hline

Warunki końcowe &
Samochód zostaje wprowadzony na stan magazynowy, otrzymuje unikalny znacznik czasu wejścia, a jego status pozwala na rozpoczęcie procesów warsztatowych. \\
\hline

Scenariusz główny &
\begin{enumerate}
\item Logistyk skanuje kod kreskowy lub wprowadza ręcznie numer VIN samochodu zjeżdżającego z lawety.
\item System odpytuje warstwę ACL integrującą się z bazą Importera/Producenta w celu weryfikacji tożsamości pojazdu.
\item System przypisuje auto do cyfrowego inwentarza salonu samochodowego i zmienia status logistyczny na \textit{Na placu (Stock)}.
\item System uruchamia wewnętrzny zegar metryki ile czasu auto spędziło na placu dla tego konkretnego samochodu.
\item System emituje zdarzenie \textit{PojazdPrzyjetyNaPlac},
\end{enumerate}
\\
\hline

Scenariusze alternatywne &
\begin{itemize}
\item [A1] Uszkodzenie w transporcie: Pracownik zauważa szkodę. System nadaje flagę \textit{Szkoda logistyczna}, wstrzymuje przyjęcie i generuje protokół szkody do ubezpieczyciela.
\item [A2] Obcy numer VIN: System odrzuca VIN jako nieprzypisany do salonu, wyrzuca błąd i blokuje przyjęcie.
\end{itemize}
\\
\hline

\end{longtable}